% =====================================================================
%  traffic_light_guide.tex
%  A complete beginner-to-viva teaching guide for the
%  YOLOv8s Traffic-Light Detection project (LISA dataset, 3 classes).
%
%  -------------------------------------------------------------------
%  HOW TO COMPILE
%  -------------------------------------------------------------------
%  This file uses the `minted` package for code highlighting, which
%  needs Python + Pygments installed (pip install Pygments) AND the
%  -shell-escape flag. Compile TWICE so the table of contents and all
%  \ref / \hyperref cross-references resolve:
%
%      pdflatex -shell-escape traffic_light_guide.tex
%      pdflatex -shell-escape traffic_light_guide.tex
%
%  (xelatex -shell-escape also works.)
%
%  -------------------------------------------------------------------
%  NO PYGMENTS?  FALLBACK TO lstlisting
%  -------------------------------------------------------------------
%  If you cannot install Pygments (e.g. Overleaf free tier limits, or a
%  bare MiKTeX install), do this:
%    1. Comment out:  \usepackage{minted}
%    2. Uncomment the listings block marked "FALLBACK" below.
%    3. Replace every  \begin{minted}{python} ... \end{minted}
%       with           \begin{lstlisting} ... \end{lstlisting}
%       (a simple find/replace).
%  Then compile normally with plain pdflatex (no -shell-escape needed).
% =====================================================================

\documentclass[11pt,a4paper]{report}

% ---------- Encoding & fonts ----------
\usepackage[utf8]{inputenc}
\usepackage[T1]{fontenc}
\usepackage{lmodern}

% ---------- Page geometry ----------
\usepackage[a4paper,margin=2.5cm]{geometry}

% ---------- Maths, tables, graphics ----------
\usepackage{amsmath,amssymb}
\usepackage{graphicx}
\usepackage{booktabs}
\usepackage{array}
\usepackage{caption}
\usepackage{float}
\usepackage{xcolor}

% ---------- Code highlighting (minted) ----------
\usepackage{minted}
\setminted{fontsize=\footnotesize, breaklines=true, breakanywhere=true,
           tabsize=4, frame=single, framesep=4pt, bgcolor=gray!6}

% ===== FALLBACK (uncomment if you cannot use minted) =================
% \usepackage{listings}
% \lstset{language=Python, basicstyle=\footnotesize\ttfamily,
%         keywordstyle=\color{blue}, commentstyle=\color{green!50!black},
%         stringstyle=\color{red!60!black}, frame=single,
%         breaklines=true, showstringspaces=false, numbers=left,
%         numberstyle=\tiny, backgroundcolor=\color{gray!6}}
% ====================================================================

% ---------- Callout boxes (tcolorbox) ----------
\usepackage[most]{tcolorbox}
\tcbuselibrary{breakable,skins}

\newtcolorbox{keyconcept}[1][]{breakable, enhanced, colback=blue!5!white,
  colframe=blue!55!black, fonttitle=\bfseries, title={Key Concept#1},
  left=2mm,right=2mm,top=1mm,bottom=1mm}

\newtcolorbox{analogy}[1][]{breakable, enhanced, colback=green!6!white,
  colframe=green!45!black, fonttitle=\bfseries, title={Real-Life Analogy#1},
  left=2mm,right=2mm,top=1mm,bottom=1mm}

\newtcolorbox{warningbox}[1][]{breakable, enhanced, colback=red!5!white,
  colframe=red!60!black, fonttitle=\bfseries, title={Warning / Common Pitfall#1},
  left=2mm,right=2mm,top=1mm,bottom=1mm}

\newtcolorbox{diagramprompt}[1][]{breakable, enhanced, colback=orange!6!white,
  colframe=orange!65!black, fonttitle=\bfseries\ttfamily,
  title={DIAGRAM PROMPT (recreate in draw.io)#1},
  left=2mm,right=2mm,top=1mm,bottom=1mm}

% ---------- Hyperlinks (load near last) ----------
\usepackage[colorlinks=true, linkcolor=blue!55!black, citecolor=green!45!black,
            urlcolor=magenta!70!black, bookmarks=true]{hyperref}

% ---------- A robust image helper ----------
% Inserts the image if it exists, otherwise a labelled placeholder box,
% so the document always compiles even if a PNG is missing.
\newcommand{\maybefig}[4]{% #1=path #2=width-fraction #3=caption #4=label
  \begin{figure}[H]\centering
    \IfFileExists{#1}{\includegraphics[width=#2\linewidth]{#1}}%
      {\fbox{\parbox[c][3.5cm][c]{0.8\linewidth}{\centering\ttfamily
        [image not found:\\ \detokenize{#1}]\\[2mm]\normalfont\small
        (compile in the project root so this PNG is on the graphics path)}}}
    \caption{#3}\label{#4}
  \end{figure}}

% Search these folders for images
\graphicspath{%
  {./}%
  {./runs/traffic_light_v1/}%
  {./Report Template/Report Template/Images/Chapter2/}%
  {./Report Template/Report Template/Images/Chapter3/}%
  {./Report Template/Report Template/Images/Chapter4/}%
}

\title{\Huge\bfseries Teaching Yourself Traffic-Light Detection\\[4mm]
  \Large A Complete Beginner-to-Viva Guide to a YOLOv8s\\
  Object-Detection Project on the LISA Dataset}
\author{Generated as a study companion for the\\
  \texttt{Traffic-Light-simplified} project notebook}
\date{\today}

\begin{document}

% ============================ TITLE ============================
\maketitle

\begin{abstract}
\noindent
This document teaches, from absolute zero, how a working traffic-light
\emph{detector} was built in the project notebook
\texttt{traffic-light-model.ipynb}. It explains the theory (what a
convolutional neural network is, what object detection means, how the
YOLOv8 architecture works, and what every evaluation metric means) and
the practice (how the data was cleaned, converted, split, balanced,
trained, evaluated, and exported). Every number quoted here is taken
directly from the project's own files: \texttt{runs/traffic\_light\_v1/results.csv},
\texttt{args.yaml}, \texttt{yolo\_dataset/data.yaml}, and the captured
notebook outputs. The final model is a fine-tuned \textbf{YOLOv8s}
network that maps the LISA dataset's $14$ raw labels down to
$3$ classes (\texttt{go}, \texttt{stop}, \texttt{warning}) and reaches a
best validation \textbf{mAP@0.5 of 0.768} at epoch~27. The last chapter is
a dedicated viva-preparation section with model answers to the questions
an examiner is most likely to ask.
\end{abstract}

\tableofcontents

% ====================================================================
\chapter{Introduction: What Are We Even Building?}
\label{ch:intro}
% ====================================================================

\section{The one-sentence summary}
We built a computer program that looks at a photo from a car's
dashboard camera and draws a coloured box around every traffic light it
sees, labelling each one as \texttt{go} (green), \texttt{stop} (red), or
\texttt{warning} (yellow/amber).

\begin{analogy}
Imagine teaching a brand-new driving student. At first they cannot even
\emph{find} the traffic lights in a busy street scene. You show them
thousands of photos, each with the lights circled and named. Eventually
they can glance at any new street and instantly point: ``red there,
green there.'' Our project does exactly this, but the ``student'' is a
neural network and the ``thousands of photos'' is the LISA dataset.
\end{analogy}

\section{Three tasks that sound similar but are not}
Beginners constantly mix these up. Pin them down now.

\begin{keyconcept}[: classification vs.\ detection vs.\ segmentation]
\begin{itemize}
  \item \textbf{Classification} answers \emph{``what is in this image?''}
        with one label for the \emph{whole} image (e.g.\ ``this is a cat'').
        It does \emph{not} say \emph{where}.
  \item \textbf{Object detection} answers \emph{``what objects are here
        and where?''} It outputs a \textbf{bounding box} (a rectangle)
        plus a class label for \emph{each} object. \textbf{This is our task.}
  \item \textbf{Segmentation} goes further and labels \emph{every pixel}
        (the exact silhouette). We do \emph{not} do this.
\end{itemize}
\end{keyconcept}

\begin{analogy}
Classification is a librarian saying ``this book is a thriller.''
Detection is the librarian saying ``there are three thrillers, and here
is the exact shelf position of each.'' Segmentation is tracing the
outline of every single book with a pen.
\end{analogy}

\section{Why traffic-light detection matters}
A self-driving or driver-assistance (ADAS, \emph{Advanced Driver
Assistance System}) car must know, many times per second, whether the
light ahead is red or green. A wrong answer is a safety failure. That is
why, as we will see in Chapter~\ref{ch:eval}, we care most about the
model never \emph{missing} a red \texttt{stop} light.

\section{The project goal, stated precisely}
Take the LISA Traffic Light dataset, simplify its $14$ fine-grained
labels into $3$ meaningful driving actions, train a YOLOv8s detector to
find and classify these lights in $1280\times960$ dashcam frames, and
produce a deployable model file. Chapters~\ref{ch:cnn}--\ref{ch:deploy}
follow the notebook cell by cell.

% ====================================================================
\chapter{CNN Foundations for Complete Beginners}
\label{ch:cnn}
% ====================================================================

\section{An image is just a grid of numbers}
\begin{keyconcept}[: pixels and tensors]
A colour image is a 3-D block of numbers called a \textbf{tensor}. Our
frames are $1280$ wide $\times\,960$ tall $\times\,3$ colour channels
(Red, Green, Blue). Each number is a pixel brightness from $0$ to $255$.
So one frame is $1280\times960\times3 = 3{,}686{,}400$ numbers. The
network only ever sees numbers --- never ``a traffic light.''
\end{keyconcept}

\begin{analogy}
Think of a mosaic made of tiny coloured tiles. From far away you see a
picture; up close it is just thousands of individual tiles. A computer
is always ``up close'': it sees the tiles (pixels), and our job is to
teach it to recognise the picture they form.
\end{analogy}

\section{Why not just use a normal neural network?}
A plain (``fully connected'') network would give every one of those
$3.6$ million numbers its own separate weight, for every neuron. That is
astronomically many parameters, and worse, it has no idea that a pixel
and its neighbour are related. It would also have to re-learn what a
traffic light looks like separately for every position in the image.

\section{Convolution: the sliding stamp}
\begin{keyconcept}[: convolution and feature maps]
A \textbf{convolution} slides a small window (a \textbf{kernel} or
\textbf{filter}, e.g.\ $3\times3$) across the image. At each position it
multiplies the pixels under it by the kernel's weights and adds them up,
producing one number. Sweeping the whole image gives a new grid called a
\textbf{feature map}. Different kernels detect different things: edges,
corners, blobs of red, round bright shapes. Stacking many convolution
layers lets later layers combine simple features (edges) into complex
ones (a traffic-light housing with a glowing lamp).
\end{keyconcept}

\begin{analogy}
A convolution kernel is a small rubber stamp. You press the same stamp
all over a page; wherever the pattern underneath matches the stamp, you
get a strong mark. Because it is the \emph{same} stamp everywhere, the
network recognises a red light whether it appears top-left or
bottom-right --- this is called \textbf{translation invariance}, and it is
why CNNs beat plain networks on images.
\end{analogy}

\begin{warningbox}
``CNN'' is not magic pattern recognition that ``just knows'' a traffic
light. It only knows what it was \emph{shown} during training. Our model
was shown small, far-away lights, so (as the notebook's own ``next
steps'' admit) it struggles with large close-up lights it never saw.
\end{warningbox}

% ====================================================================
\chapter{From Classification to Detection: How YOLO Works}
\label{ch:yolo}
% ====================================================================

\section{Bounding boxes}
A \textbf{bounding box} is the tightest rectangle that encloses an
object. It is described by four numbers. The dataset stores two opposite
\emph{corners} in pixels; YOLO instead wants the box \emph{centre} plus
its width and height, each divided by the image size so the values sit
between $0$ and $1$ (we convert between the two in
Chapter~\ref{ch:dataprep}).

\section{IoU: measuring ``how good is this box?''}
\begin{keyconcept}[: Intersection over Union (IoU)]
\textbf{IoU} compares a predicted box with the true box:
\[
\mathrm{IoU} = \frac{\text{area of overlap}}{\text{area of union}}.
\]
$\mathrm{IoU}=1$ means a perfect match; $0$ means no overlap. A
prediction usually ``counts'' as correct only if its IoU with a real box
exceeds a threshold (commonly $0.5$). IoU also drives \textbf{NMS}
(Section~\ref{sec:nms}).
\end{keyconcept}

\begin{analogy}
Two people describe where a parked car is by drawing chalk rectangles on
the road. IoU is how much their two rectangles overlap. Lots of
overlap = they basically agree.
\end{analogy}

\section{One-stage vs.\ two-stage, and why we chose YOLO}
\begin{keyconcept}[: one-stage detectors]
\textbf{Two-stage} detectors (e.g.\ Faster R-CNN) first propose
candidate regions, then classify each --- accurate but slower.
\textbf{One-stage} detectors like \textbf{YOLO} (``You Only Look Once'')
predict all boxes and classes in a single pass over the image --- fast
enough for real-time video, which a car needs. YOLOv8 is
\textbf{anchor-free}: instead of comparing against a fixed set of
pre-defined ``anchor'' box shapes, it directly predicts each box's
geometry, which simplifies training.
\end{keyconcept}

\section{The YOLOv8 architecture: backbone, neck, head}
\label{sec:arch}
\begin{keyconcept}[: the three parts of YOLOv8]
\begin{itemize}
  \item \textbf{Backbone} --- a deep CNN that turns the raw image into
        feature maps at several resolutions. It extracts ``what is
        where, roughly.''
  \item \textbf{Neck} (a feature-pyramid, FPN/PAN style) --- mixes the
        fine, high-resolution maps (good for \emph{small} objects) with
        the coarse, deep maps (good for \emph{large} objects and
        context). Crucial here because traffic lights are tiny.
  \item \textbf{Head} --- the output layer that, for each location,
        predicts the box geometry and the class scores
        (\texttt{go}/\texttt{stop}/\texttt{warning}).
\end{itemize}
Our specific model, \textbf{YOLOv8s} (``small''), has
\textbf{11{,}126{,}745 parameters} and runs at $28.4$ GFLOPs
(from the notebook's model summary).
\end{keyconcept}

\begin{analogy}
A factory line: the \textbf{backbone} is workers who inspect the raw
material and note its features; the \textbf{neck} is a supervisor who
combines notes from close-up and wide-angle inspectors; the
\textbf{head} is the stamper who writes the final verdict ``red light,
here, this big.''
\end{analogy}

\begin{diagramprompt}
Draw three large boxes left-to-right joined by arrows:
[INPUT IMAGE 1280x1280x3] -> [BACKBONE: stack of Conv + C2f blocks,
shrinking H,W while growing channels; show 3 side outputs P3,P4,P5] ->
[NECK: FPN/PAN -- two columns of arrows, one going up (upsample + concat),
one going down (downsample + concat), merging P3/P4/P5] ->
[HEAD: 3 detection branches, each outputting (box: x,y,w,h) + (3 class
scores)]. Label P3 "fine map / small objects", P5 "coarse map / large
objects". Add a final arrow to [NMS] -> [FINAL BOXES].
\end{diagramprompt}

\maybefig{runs/traffic_light_v1/training_curves.png}{0.95}{The six training curves YOLOv8
saved for this run. We dissect these in Chapter~\ref{ch:train}.}{fig:curves-preview}

% ====================================================================
\chapter{The Dataset: LISA Traffic Lights}
\label{ch:data}
% ====================================================================

\section{What we were given}
The source is the \textbf{LISA Traffic Light Dataset} in Supervisely
JSON format. From the notebook's verification output (CELL~2):
\begin{itemize}
  \item \textbf{Train:} $20{,}535$ images, each with a matching
        \texttt{.jpg.json} annotation file.
  \item \textbf{Test:} $22{,}481$ images.
  \item \textbf{Total:} $43{,}016$ images. Every frame is
        $1280\times960$ pixels.
  \item \texttt{meta.json} declares $14$ classes.
  \item Frames are grouped into \textbf{sequences} (short video clips
        like \texttt{dayClip1}, \texttt{daySequence1}); consecutive
        frames look almost identical. Remember this --- it causes the
        ``data leakage'' problem in Chapter~\ref{ch:hardproblems}.
\end{itemize}

\section{The annotation format}
Each object is a \texttt{rectangle} with two corner points in pixels.
Here is the real first training file (CELL~3), trimmed:

\begin{minted}{json}
{
  "size": { "height": 960, "width": 1280 },
  "tags": [ {"name": "sequence", "value": "Clip1"}, {"name": "day"} ],
  "objects": [
    { "classTitle": "go traffic light",
      "points": { "exterior": [ [698, 333], [710, 358] ] } }
  ]
}
\end{minted}

\noindent
So this light's box runs from corner $(698,333)$ to $(710,358)$ --- only
$12\times25$ pixels, about $0.024\%$ of the image. Traffic lights are
\emph{tiny}.

\section{Why collapse 14 classes into 3?}
The raw labels separate things a \emph{driver} does not care about:
``stop'' vs.\ ``stop left'' vs.\ ``stop left traffic light'' all mean
\emph{the car must stop}. Splitting hairs would (a) starve each tiny
class of training examples and (b) add no driving value. So CELL~5 maps
everything onto three \emph{actions}:

\begin{keyconcept}[: the 14$\to$3 mapping]
\texttt{go}, \texttt{go left}, \texttt{go forward} ($+$ ``traffic light''
variants) $\rightarrow$ \textbf{go}.\quad
\texttt{stop}, \texttt{stop left} ($+$ variants) $\rightarrow$
\textbf{stop}.\quad
\texttt{warning}, \texttt{warning left} ($+$ variants) $\rightarrow$
\textbf{warning}.\\
Numeric IDs (YOLO needs numbers, not words):
\texttt{go}$=0$, \texttt{stop}$=1$, \texttt{warning}$=2$.
\end{keyconcept}

\section{Exploratory data analysis (EDA): three findings}
\subsection{Finding 1 --- severe class imbalance}
After mapping, the $104{,}534$ training boxes split as in
Table~\ref{tab:classdist}.

\begin{table}[H]\centering
\caption{Class distribution after the 14$\to$3 mapping (CELL~5/6 output).}
\label{tab:classdist}
\begin{tabular}{lrr}
\toprule
Class & Boxes & Share \\
\midrule
\texttt{stop}    & $52{,}450$ & $50.2\%$ \\
\texttt{go}      & $48{,}892$ & $46.8\%$ \\
\texttt{warning} & $3{,}192$  & $3.1\%$  \\
\midrule
Total & $104{,}534$ & $100\%$ \\
\bottomrule
\end{tabular}
\end{table}

\noindent
\texttt{warning} is starved: the ratio is about
\textbf{16:1} (stop\,:\,warning). We fix this in
Chapter~\ref{ch:hardproblems}.

\maybefig{class_distribution.png}{0.95}{Left: raw 14-class counts.
Right: the consolidated 3 classes, with the under-represented
\texttt{warning} bar flagged in red.}{fig:classdist}

\subsection{Finding 2 --- these are small objects}
From CELL~7: the median box side is only $\approx 18$~px and covers
$0.0263\%$ of the frame. Using the COCO size convention,
\textbf{$68.6\%$ of boxes are ``Tiny'' ($<32$\,px)} and $31.3\%$ are
``Small'' --- $99.9\%$ together. This single fact justifies two later
choices: the high $1280$ input resolution and the YOLOv8\textbf{s} (not
nano) model.

\maybefig{bbox_analysis.png}{0.95}{Box-size analysis: width-vs-height
scatter, relative-width histogram, and the COCO size-category pie. The
verdict is ``small-object detection problem.''}{fig:bbox}

\subsection{Finding 3 --- day/night and sequences}
Of images with objects, $12{,}872$ are daytime and $5{,}266$ nighttime;
the data spans $18$ sequences. Night is under-represented (a noted
weakness).

\maybefig{class_samples.png}{0.9}{Sample frames per class. The inner box
is the true annotation; the outer box is a zoom guide because the real
lights are too small to see otherwise.}{fig:samples}

% ====================================================================
\chapter{Data Preparation: From JSON to YOLO}
\label{ch:dataprep}
% ====================================================================

\section{The folder layout YOLO expects}
YOLO finds a label by taking the image path, swapping the word
\texttt{images} for \texttt{labels}, and \texttt{.jpg} for \texttt{.txt}
(CELL~9):
\begin{minted}{text}
yolo_dataset/
  data.yaml
  images/{train,val,test}/   <- the .jpg frames
  labels/{train,val,test}/   <- one .txt per image (the boxes)
\end{minted}

\section{Converting one box (the maths)}
\begin{keyconcept}[: normalised YOLO coordinates]
Given corner pixels $(x_1,y_1)$ and $(x_2,y_2)$ in an image
$W\times H$, YOLO wants
\[
x_c=\frac{x_1+x_2}{2W},\;\; y_c=\frac{y_1+y_2}{2H},\;\;
w=\frac{x_2-x_1}{W},\;\; h=\frac{y_2-y_1}{H}.
\]
Dividing by $W,H$ makes the numbers \emph{resolution-independent}
(always $0$--$1$), so the same label still works if the image is resized.
\end{keyconcept}

\paragraph{Worked example (the real first box).}
Corners $(698,333)\to(710,358)$ in $1280\times960$:
\[
x_c=\tfrac{698+710}{2\cdot1280}=0.5500,\;
y_c=\tfrac{333+358}{2\cdot960}=0.3599,\;
w=\tfrac{12}{1280}=0.0094,\;
h=\tfrac{25}{960}=0.0260.
\]
The label line is therefore \texttt{0 0.550000 0.359896 0.009375 0.026042}
(class~$0=$\,\texttt{go}) --- exactly what the notebook printed.

\begin{minted}{python}
# CELL 10 (core of the converter) -- see the notebook for the full loop
x_centre = (x1 + x2) / 2 / img_w
y_centre = (y1 + y2) / 2 / img_h
width    = (x2 - x1)     / img_w
height   = (y2 - y1)     / img_h
line = f"{class_id} {x_centre:.6f} {y_centre:.6f} {width:.6f} {height:.6f}"
\end{minted}

\begin{diagramprompt}
Two side-by-side panels. LEFT "JSON / pixel space": a 1280x960 rectangle
with a tiny box; annotate corners (698,333) top-left and (710,358)
bottom-right. RIGHT "YOLO / normalised space": the same rectangle scaled
to 1x1; annotate centre dot (0.550, 0.360), width 0.0094, height 0.0260.
A big arrow labelled "divide by W,H" points LEFT->RIGHT.
\end{diagramprompt}

\section{The \texttt{data.yaml} file}
This tiny file is what you actually hand to YOLO. The project's real one:
\begin{minted}{yaml}
path: .../yolo_dataset
train: images/train
val:   images/val
test:  images/test
nc: 3
names: [go, stop, warning]   # index 0, 1, 2 -- order MUST match the IDs
\end{minted}

\begin{warningbox}
The order of \texttt{names} \emph{is} the class ID. If you wrote
\texttt{[stop, go, warning]} here, every label file (which uses $0$ for
\texttt{go}) would silently be mis-named and your model would learn red
lights as ``go.'' CELL~13 explicitly validates this order.
\end{warningbox}

\section{Trust, but verify: auditing and visualising}
Two safety nets before training:
\begin{itemize}
  \item \textbf{Audit (CELL~14):} every \texttt{.txt} line must have
        exactly $5$ values, a class ID in $\{0,1,2\}$, and coordinates in
        range. The project reported \textbf{0 errors} across all splits.
  \item \textbf{Visual check (CELL~15):} re-draw boxes \emph{from the
        generated \texttt{.txt} files} (not the JSON). If they still land
        on the lights, the conversion was correct.
\end{itemize}

\maybefig{yolo_label_check.png}{0.92}{Boxes drawn from the converted
\texttt{.txt} labels. Because they still sit on the lights, the JSON
$\to$ YOLO conversion is verified.}{fig:labelcheck}

% ====================================================================
\chapter{The Two Hard Problems (and Their Fixes)}
\label{ch:hardproblems}
% ====================================================================
These two decisions are the most ``defendable'' engineering choices in
the project --- examiners love them.

\section{Problem 1: data leakage from near-identical frames}
\begin{keyconcept}[: data leakage]
\textbf{Data leakage} is when information from the test/validation data
sneaks into training, giving a falsely high score. Here the danger is
subtle: consecutive frames in a clip are almost the same picture. If
\texttt{dayClip1--00000} went to \emph{train} and
\texttt{dayClip1--00001} went to \emph{val}, the model would be
``tested'' on a frame it has practically already seen.
\end{keyconcept}

\begin{analogy}
It is like giving a student the exam questions as homework the night
before. They score $99\%$, but you have learned nothing about whether
they actually understand the subject.
\end{analogy}

\paragraph{The fix --- sequence-aware split (CELL~11).} Split by
\emph{whole sequence}, never by frame. Of $18$ sequences, $15$ went to
train ($15{,}658$ images) and $3$ to validation ($4{,}877$ images). The
code then asserts no sequence appears in both sets --- the
\emph{leakage check PASSED}.

\begin{diagramprompt}
Draw a row of film-strip icons grouped into 18 labelled clips. Colour 15
clips blue ("train") and 3 green ("val"). Show a red crossed-out arrow
between a blue and a green clip labelled "NO shared sequence -> no
leakage". Contrast with a faded "WRONG" inset where individual frames of
ONE clip are split across train and val.
\end{diagramprompt}

\begin{warningbox}
Because there are only $\sim$18 sequences and \texttt{warning}-dominant
clips are rare, a \emph{stratified} split failed, so the code falls back
to a plain random split of sequences. The trade-off: the class balance of
train vs.\ val is not perfectly controlled --- which partly explains why
per-class validation numbers vary a lot.
\end{warningbox}

\section{Problem 2: class imbalance}
With \texttt{warning} at only $\sim$3\% of boxes, a lazy model could
ignore yellow lights entirely and still score well on raw accuracy.

\begin{keyconcept}[: oversampling]
\textbf{Oversampling} means showing the rare class \emph{more often}
during training. CELL~12 finds the $404$ \texttt{warning}-dominant train
images ($2.6\%$) and creates extra copies (via symlinks) so that
\texttt{warning} rises to a target. Using a $5\times$ factor lifts the
warning share to about \textbf{11.7\%} of the train set.
\end{keyconcept}

\begin{analogy}
If a student keeps failing the few questions on one rare topic, you make
them practise that topic five times as much. They do not get new facts,
just more exposure, until it sticks.
\end{analogy}

\begin{warningbox}
Oversampling shows the \emph{same} images repeatedly, which risks
\textbf{overfitting} to those exact frames. The project mitigates this by
relying on YOLO's \textbf{mosaic} and \textbf{mixup} augmentation, which
recombine and blend images so each ``copy'' is seen in a fresh context.
\end{warningbox}

% ====================================================================
\chapter{Training the Model}
\label{ch:train}
% ====================================================================

\section{The training command}
Training is a few lines (CELL~16); the heavy lifting is Ultralytics'.

\begin{minted}{python}
from ultralytics import YOLO
model = YOLO('yolov8s.pt')          # start from COCO-pretrained weights
results = model.train(
    data='yolo_dataset/data.yaml',
    epochs=50, imgsz=1280, batch=16, device=0,
    mosaic=1.0, mixup=0.1, fliplr=0.5,           # augmentation
    lr0=0.01, lrf=0.01, warmup_epochs=3,         # optimiser schedule
    patience=20, seed=42, name='traffic_light_v1')
\end{minted}

\begin{keyconcept}[: the key hyperparameters, in plain words]
\begin{itemize}
  \item \textbf{epoch} --- one full pass over the training set. We allowed
        up to $50$.
  \item \textbf{batch} ($=16$) --- how many images are processed together
        before the weights update once.
  \item \textbf{imgsz} ($=1280$) --- images are fed at $1280$ px so the
        tiny lights survive; lower res would erase them.
  \item \textbf{transfer learning} --- we start from \texttt{yolov8s.pt},
        already trained on COCO (80 classes), and \emph{fine-tune} it on
        our 3 classes. The network already knows edges/shapes, so it
        learns fast.
  \item \textbf{patience} ($=20$) --- ``early stopping'': if validation
        mAP does not improve for $20$ epochs, stop and keep the best
        weights.
  \item \textbf{seed} ($=42$) --- fixes randomness so the run is
        reproducible.
\end{itemize}
\end{keyconcept}

\begin{analogy}[: transfer learning]
You do not teach a new employee to read before teaching them your filing
system. They already read (COCO pre-training); you just teach the
specifics of your office (the 3 traffic-light classes). Far faster than
raising someone from birth.
\end{analogy}

\section{What the three losses mean}
YOLOv8 minimises a sum of three losses (their weights are in
\texttt{args.yaml}: \texttt{box}$=7.5$, \texttt{cls}$=0.5$,
\texttt{dfl}$=1.5$):
\begin{itemize}
  \item \textbf{box\_loss} --- how wrong the box \emph{position/size} is.
  \item \textbf{cls\_loss} --- how wrong the \emph{class}
        (go/stop/warning) is.
  \item \textbf{dfl\_loss} (Distribution Focal Loss) --- sharpens the
        \emph{exact edges} of the box.
\end{itemize}
A good run shows all three falling on both train and val curves.

\section{Reading the training log}
The real progression from \texttt{results.csv}:

\begin{table}[H]\centering
\caption{Validation mAP@0.5 at key epochs (from \texttt{results.csv}).}
\label{tab:milestones}
\begin{tabular}{cll}
\toprule
Epoch & mAP@0.5 & Meaning \\
\midrule
$1$  & $0.664$ & barely started \\
$10$ & $0.706$ & rapid learning \\
$\mathbf{27}$ & $\mathbf{0.768}$ & \textbf{best --- saved as best.pt} \\
$47$ & $0.667$ & early-stopped (no gain for 20 epochs) \\
\bottomrule
\end{tabular}
\end{table}

\begin{keyconcept}[: why \texttt{best.pt} is epoch 27, not epoch 47]
YOLO saves two files: \texttt{last.pt} (most recent) and
\texttt{best.pt} (the epoch with the highest validation mAP). Validation
performance peaked at epoch~27 ($\text{mAP@0.5}=0.768$,
$\text{mAP@0.5:0.95}=0.379$) and then drifted down --- a sign of mild
\textbf{overfitting}. Early stopping triggered at epoch~47. We deploy
\texttt{best.pt} (epoch~27), never the last one. Total training time
$\approx$ 2\,h\,14\,min on an RTX~5090.
\end{keyconcept}

\maybefig{runs/traffic_light_v1/training_curves.png}{0.95}{Top row: the three losses (train
solid, val dashed) all trending down. Bottom row: mAP@0.5, mAP@0.5:0.95,
and precision/recall. The red dotted line marks the best epoch.}{fig:curves}

\begin{warningbox}
Notice the validation losses are \emph{noisier} and flatten higher than
the training losses --- normal, but the widening gap after epoch~27 is the
overfitting signal that justifies keeping the early-stopped best weights.
\end{warningbox}

% ====================================================================
\chapter{Evaluation: Did It Actually Work?}
\label{ch:eval}
% ====================================================================

\section{The four metrics, demystified}
\begin{keyconcept}[: precision, recall, mAP]
For a class:
\begin{itemize}
  \item \textbf{Precision} $=\dfrac{TP}{TP+FP}$ --- of the boxes the model
        \emph{claimed}, how many were right. Low precision $=$ ``cries
        wolf.''
  \item \textbf{Recall} $=\dfrac{TP}{TP+FN}$ --- of the real lights, how
        many it \emph{found}. Low recall $=$ ``misses things.''
  \item \textbf{mAP@0.5} --- mean Average Precision with the IoU bar at
        $0.5$; the headline accuracy score, averaged over classes.
  \item \textbf{mAP@0.5:0.95} --- the average of mAP at IoU thresholds
        $0.5,0.55,\dots,0.95$. Much stricter (the box must be
        \emph{very} tight), so it is always lower.
\end{itemize}
($TP$ = true positive, $FP$ = false positive, $FN$ = missed object.)
\end{keyconcept}

\begin{analogy}
A fisherman's net. \textbf{Precision}: of everything in the net, how much
is actually fish (not boots)? \textbf{Recall}: of all the fish in the
lake, how many did the net catch? You usually trade one for the other.
\end{analogy}

\section{Per-class results}
Table~\ref{tab:results} is the project's own per-class breakdown
(CELL~23, validation set, \texttt{best.pt}). The overall row matches the
epoch-27 figures in \texttt{results.csv}.

\begin{table}[H]\centering
\caption{Per-class metrics for \texttt{best.pt} (validation set).}
\label{tab:results}
\begin{tabular}{lrrr}
\toprule
Class & Precision & Recall & mAP@0.5 \\
\midrule
\texttt{go}      & $0.942$ & $0.559$ & $0.645$ \\
\texttt{stop}    & $0.790$ & $0.959$ & $0.816$ \\
\texttt{warning} & $0.891$ & $0.821$ & $0.843$ \\
\midrule
\textbf{ALL}     & $0.874$ & $0.780$ & $0.768$ \\
\bottomrule
\end{tabular}
\end{table}

\paragraph{How to read this like an examiner.}
\begin{itemize}
  \item \texttt{stop} has \textbf{recall $0.959$} --- the model almost
        never \emph{misses} a red light. For safety, this is the single
        most important number, and it is excellent.
  \item \texttt{go} has high precision ($0.942$) but \textbf{low recall
        ($0.559$)} --- when it says ``green'' it is usually right, but it
        \emph{misses} many green lights. The notebook flags this as the
        top weakness to fix.
  \item \texttt{warning} did \emph{well} ($\text{mAP}=0.843$,
        recall $0.821$) --- direct evidence that the oversampling in
        Chapter~\ref{ch:hardproblems} worked.
\end{itemize}

\maybefig{runs/traffic_light_v1/test_metrics.png}{0.8}{Per-class precision, recall and mAP@0.5
as bars. The dashed line at $0.8$ is a quality reference.}{fig:testmetrics}

\section{The confusion matrix}
\begin{keyconcept}[: confusion matrix]
A grid where each column is a \emph{true} class and each row is the
\emph{predicted} class (plus a ``background'' row/column for
misses/false alarms). Bright cells on the diagonal $=$ correct. Bright
\emph{off-diagonal} cells $=$ the model confusing one class for another;
the \texttt{background} column shows real lights that were \emph{missed}.
\end{keyconcept}

\maybefig{runs/traffic_light_v1/confusion_matrix_normalized.png}{0.7}{Normalised confusion
matrix for \texttt{best.pt}. A strong diagonal means the three colours
are rarely confused with each other; the bulk of error is the
\texttt{background} row --- i.e.\ small lights missed entirely, consistent
with \texttt{go}'s low recall.}{fig:cm}

\section{Qualitative check --- look at predictions}
Numbers can lie; pictures rarely do. CELL~20 runs the model on random
test frames at a confidence threshold of $0.30$.

\maybefig{runs/traffic_light_v1/sample_inference.png}{0.95}{Detections on test frames
(confidence $\ge 0.30$). Each box shows the predicted class and
confidence.}{fig:inference}

% ====================================================================
\chapter{Inference and Deployment}
\label{ch:deploy}
% ====================================================================

\section{Confidence threshold and NMS}
\label{sec:nms}
\begin{keyconcept}[: confidence threshold]
For every candidate box the model outputs a \textbf{confidence}
($0$--$1$). We keep only boxes above a threshold --- the project uses
\textbf{$0.30$}. Higher $=$ fewer, surer boxes (precision up, recall
down); lower $=$ more boxes (recall up, more false alarms).
\end{keyconcept}

\begin{keyconcept}[: Non-Maximum Suppression (NMS)]
The model often fires several overlapping boxes on one light.
\textbf{NMS} keeps the highest-confidence box and deletes others that
overlap it by more than an IoU threshold (here $0.7$, from
\texttt{args.yaml}). This is why you see one clean box per light, not ten.
\end{keyconcept}

\begin{analogy}
Ten witnesses describe the same parked car. NMS keeps the most confident
witness's account and discards the nine near-duplicates, so you record
\emph{one} car, not ten.
\end{analogy}

\section{Single image, and video from frames}
CELL~20 does single images; CELL~21 stitches sequential frames into an
annotated video. On \texttt{daySequence1} ($300$ frames $=10$\,s at
$30$\,fps) the model produced $614$ detections, averaging $2.0$ per
frame. CELL~21b runs on \emph{your own} video --- a real $1920\times1080$
clip of $5{,}403$ frames processed at $\approx 92$ fps with $10{,}686$
detections.

\section{ONNX export}
\begin{keyconcept}[: what ONNX is]
\textbf{ONNX} (Open Neural Network Exchange) is a universal model file
format. Exporting \texttt{best.pt} (PyTorch) to \texttt{best.onnx} lets
the model run \emph{without} Python or PyTorch --- e.g.\ on a phone, in a
browser, or in C++ --- via lightweight ``ONNX Runtime.''
\end{keyconcept}

\begin{minted}{python}
# CELL 22
model.export(format='onnx', imgsz=1280, dynamic=False,
             opset=17, simplify=True)
\end{minted}

\noindent
Result (CELL~22 output): \texttt{best.pt} ($22.6$\,MB) $\rightarrow$
\texttt{best.onnx} ($45.3$\,MB), input named \texttt{images}
$(1,3,1280,1280)$, output \texttt{output0}. Validation passed.

\section{Three deployment routes (CELL~22 notes)}
\begin{itemize}
  \item \textbf{Android app} --- bundle \texttt{best.onnx} as an asset,
        run it with ONNX Runtime Mobile on camera frames. Works offline,
        no internet, all on-device.
  \item \textbf{Website, in-browser} --- ONNX Runtime Web (WASM) runs the
        model in the visitor's browser. No server, privacy-friendly,
        but limited by the user's device.
  \item \textbf{Website, server backend} --- the browser sends the image
        to an API; the server runs the model and returns boxes. Most
        powerful hardware, but needs a server and a network round-trip.
\end{itemize}

\begin{warningbox}
The model was trained at \texttt{imgsz=1280}, which is heavy for phones.
The notebook's own advice is to \emph{re-export at \texttt{imgsz=640}}
for mobile/web --- roughly $4\times$ faster --- accepting a small accuracy
drop on the smallest lights.
\end{warningbox}

\section{Honest limitations and next steps}
Straight from the project summary, the known weaknesses are: misses
\emph{large} close-up lights (never seen in LISA); \texttt{go} recall is
only $0.559$; night data is under-represented; and $1280$ px is slow for
edge devices. Sensible improvements: add your own annotated frames, try
YOLOv8m, add night augmentation, and fine-tune on domain video.

% ====================================================================
\chapter{Viva Preparation}
\label{ch:viva}
% ====================================================================
Answer out loud. Each answer below is deliberately short enough to say in
$20$--$40$ seconds.

\section{Theme A --- CNN \& detection basics}
\begin{enumerate}
  \item \textbf{What is the difference between classification and
        detection?} Classification gives one label for the whole image;
        detection gives a box \emph{and} a label for every object. We do
        detection.
  \item \textbf{What is a convolution?} A small kernel slid across the
        image computing weighted sums, producing a feature map; it
        detects local patterns and is reused at every position
        (translation invariance).
  \item \textbf{Why a CNN and not a plain neural network?} Plain
        networks ignore spatial structure and need far too many
        parameters; CNNs share weights and exploit locality.
  \item \textbf{What is a bounding box and how is it stored here?} The
        tightest rectangle around an object; stored in YOLO as
        normalised $(x_c,y_c,w,h)$ in $[0,1]$.
  \item \textbf{What is IoU?} Intersection-over-union of predicted and
        true boxes; the overlap quality measure used for matching and NMS.
\end{enumerate}

\section{Theme B --- why YOLO / the architecture}
\begin{enumerate}\setcounter{enumi}{5}
  \item \textbf{Why YOLO over Faster R-CNN?} YOLO is one-stage and
        real-time --- essential for a moving car --- with competitive
        accuracy.
  \item \textbf{What does ``anchor-free'' mean?} YOLOv8 predicts box
        geometry directly instead of refining pre-set anchor shapes;
        simpler, fewer hyperparameters.
  \item \textbf{Backbone, neck, head?} Backbone extracts multi-scale
        features; neck (FPN/PAN) fuses fine and coarse maps; head outputs
        boxes $+$ class scores.
  \item \textbf{Why YOLOv8\emph{s} and not nano or medium?} EDA showed
        $99.9\%$ tiny/small objects; ``s'' detects small objects far
        better than nano while fitting our GPU, where medium was
        unnecessary.
  \item \textbf{How many parameters?} About $11.1$ million
        ($11{,}126{,}745$), $28.4$ GFLOPs.
\end{enumerate}

\section{Theme C --- data, leakage, imbalance}
\begin{enumerate}\setcounter{enumi}{10}
  \item \textbf{Why map 14 classes to 3?} The extra labels are
        driving-irrelevant and would starve each class; 3 actions
        (go/stop/warning) are what a car needs.
  \item \textbf{What is data leakage and how did you prevent it?}
        Train/test overlap that inflates scores. Frames in a clip are
        near-identical, so we split by whole \emph{sequence} (15 train /
        3 val) and verified no sequence is shared.
  \item \textbf{What is class imbalance and your fix?} \texttt{warning}
        was $\sim$3\%; we oversampled warning-dominant images $5\times$ to
        $\sim$11.7\% of the train set.
  \item \textbf{Doesn't oversampling cause overfitting?} It can; mosaic
        and mixup augmentation place each repeated image in fresh
        contexts to counter that.
  \item \textbf{Why did stratified splitting fail?} Only $\sim$18
        sequences exist and warning-dominant ones are rare, so some
        strata had a single member; the code falls back to a plain split.
\end{enumerate}

\section{Theme D --- training \& metrics}
\begin{enumerate}\setcounter{enumi}{15}
  \item \textbf{What are the three losses?} box (position/size), cls
        (class), dfl (edge sharpness).
  \item \textbf{What is transfer learning here?} Start from
        COCO-pretrained \texttt{yolov8s.pt} and fine-tune on 3 classes;
        faster and better than training from scratch.
  \item \textbf{Why is \texttt{best.pt} from epoch 27?} Validation mAP
        peaked there ($0.768$) then declined (overfitting); early stopping
        fired at epoch~47 but we keep the best.
  \item \textbf{Define precision, recall, mAP@0.5, mAP@0.5:0.95.} See
        Chapter~\ref{ch:eval}; recall $=$ fraction of real lights found,
        precision $=$ fraction of claims that are correct, mAP $=$ mean
        average precision at the stated IoU bar(s).
  \item \textbf{Your overall scores?} mAP@0.5 $=0.768$, mAP@0.5:0.95
        $=0.379$, precision $0.874$, recall $0.780$ (validation, epoch 27).
\end{enumerate}

\section{Theme E --- the confusion matrix \& errors}
\begin{enumerate}\setcounter{enumi}{20}
  \item \textbf{Which class is best and why does it matter?} \texttt{stop}
        recall $=0.959$ --- safety-critical, and excellent: red lights are
        rarely missed.
  \item \textbf{Which class is weakest?} \texttt{go} recall $=0.559$ ---
        many green lights missed; high precision means few false greens.
  \item \textbf{Where does most error come from in the confusion
        matrix?} The \texttt{background} row --- small/far lights missed
        entirely --- not go/stop confusion.
  \item \textbf{Is mAP@0.5:0.95 low ($0.379$) a problem?} It is a
        \emph{strict} metric requiring very tight boxes on tiny objects;
        for ``which colour, roughly where'' the $0.768$ mAP@0.5 is the
        relevant number.
\end{enumerate}

\section{Theme F --- deployment, ethics, improvements}
\begin{enumerate}\setcounter{enumi}{24}
  \item \textbf{What is ONNX and why export?} A portable model format to
        run without PyTorch --- on phones, browsers, or C++ via ONNX
        Runtime.
  \item \textbf{What is NMS and the thresholds used?} Removes duplicate
        overlapping boxes (IoU $0.7$); we keep detections with confidence
        $\ge 0.30$.
  \item \textbf{Deployment options?} On-device Android (offline),
        in-browser WASM (no server), or server API (most power).
  \item \textbf{Safety/ethics concerns?} A missed red light is dangerous;
        the model must be a \emph{driver aid}, not sole authority, and
        needs validation in rain, glare, and night before any real use.
  \item \textbf{Biggest real-world weakness?} It only saw small, mostly
        daytime LISA lights; it under-performs on close-up lights and at
        night.
  \item \textbf{Three concrete improvements?} Add annotated domain
        frames; try YOLOv8m; add night-specific augmentation (and
        re-export at $640$ for speed).
  \item \textbf{If you had one more week?} Raise \texttt{go} recall:
        oversample go-dominant clips, tune augmentation, and lower the
        confidence threshold while watching precision.
\end{enumerate}

\section{One-page cheat-sheet (memorise one line per topic)}
\begin{keyconcept}[: the essentials]
\begin{itemize}
  \item \textbf{Task:} object detection (box $+$ label), not just
        classification.
  \item \textbf{Model:} YOLOv8s, one-stage, anchor-free, $11.1$M params,
        fine-tuned from COCO.
  \item \textbf{Data:} LISA, $20{,}535$ train / $22{,}481$ test images,
        $14\to3$ classes (go/stop/warning), $99.9\%$ tiny objects.
  \item \textbf{Key fix 1:} sequence-aware split $\Rightarrow$ no data
        leakage.
  \item \textbf{Key fix 2:} $5\times$ oversample \texttt{warning}
        ($3\%\to12\%$).
  \item \textbf{Training:} imgsz $1280$, batch $16$, up to $50$ epochs,
        best at epoch $27$, early-stopped at $47$.
  \item \textbf{Result:} mAP@0.5 $=0.768$; \texttt{stop} recall $0.959$
        (great), \texttt{go} recall $0.559$ (weak).
  \item \textbf{Deploy:} export to ONNX; phone / browser / server.
\end{itemize}
\end{keyconcept}

% ====================================================================
\chapter*{Glossary}
\addcontentsline{toc}{chapter}{Glossary}
% ====================================================================
\begin{description}
  \item[ADAS] Advanced Driver Assistance System --- safety tech that
        helps a human driver.
  \item[Anchor-free] Predicting boxes directly rather than refining
        preset anchor shapes.
  \item[Backbone/Neck/Head] Feature extractor / multi-scale fuser /
        output predictor in a detector.
  \item[Batch] Number of images processed before one weight update
        ($16$ here).
  \item[Bounding box] Rectangle around an object; YOLO stores it as
        normalised $(x_c,y_c,w,h)$.
  \item[CNN] Convolutional Neural Network --- a network built from
        convolution layers.
  \item[Confidence threshold] Minimum score to keep a detection
        ($0.30$ here).
  \item[Data leakage] Train/test contamination that inflates scores.
  \item[Epoch] One full pass over the training data.
  \item[IoU] Intersection over Union --- box overlap measure.
  \item[mAP@0.5 / mAP@0.5:0.95] Mean Average Precision at one / a range
        of IoU thresholds.
  \item[NMS] Non-Maximum Suppression --- removes duplicate overlapping
        boxes.
  \item[ONNX] Portable model format for framework-free inference.
  \item[Oversampling] Repeating rare-class examples to fight imbalance.
  \item[Precision / Recall] Fraction of claims correct / fraction of real
        objects found.
  \item[Transfer learning] Fine-tuning a pretrained network on a new
        task.
  \item[YOLO] ``You Only Look Once'' --- a family of one-stage detectors.
\end{description}

% ====================================================================
\begin{thebibliography}{9}
\addcontentsline{toc}{chapter}{Bibliography}
% ====================================================================
\bibitem{ultralytics}
Ultralytics, \emph{YOLOv8 Documentation}.
\url{https://docs.ultralytics.com}.

\bibitem{yolo}
J.~Redmon, S.~Divvala, R.~Girshick, and A.~Farhadi,
``You Only Look Once: Unified, Real-Time Object Detection,''
\emph{CVPR}, 2016. \url{https://arxiv.org/abs/1506.02640}.

\bibitem{lisa}
M.~P.~Philipsen \emph{et al.}, ``The LISA Traffic Light Dataset,''
University of California, San Diego.
\url{https://www.kaggle.com/datasets/mbornoe/lisa-traffic-light-dataset}.

\bibitem{coco}
T.-Y.~Lin \emph{et al.}, ``Microsoft COCO: Common Objects in Context,''
\emph{ECCV}, 2014. \url{https://cocodataset.org}.

\bibitem{onnx}
ONNX, \emph{Open Neural Network Exchange}. \url{https://onnx.ai}.
\end{thebibliography}

\end{document}
